\magnification=\magstep1
\input notesmac.tex
\nopagenumbers


\centerline{\titlefont CZ3106\ \  Symbolic Computing, Tutorial 6}
\bigskip
\bigskip
\line{A/P~Wang Jian-Sheng \hfill week 7-12 Apr 03}
\bigskip
\bigskip


\bigskip
\problem Solve the following systems of linear equations (by hand) with
Chinese remainder theorem (Garner's algorithm) by solving the
equations in Z$_p$ for $p=3, 5, 7, 11$ first, and then reconstruct the
solution in Z.  Of course the answer may be in Q. 
$$ 2x + 4y - 2z = 1,$$
$$ 2x + y + z = 3, $$
$$ 5x - 3y + 7z = 0.$$
An example of similar calculation is given on page 182--183, ``Algorithm
for Computer Algebra'', by Geddes et al.

\bigskip
\problem Calculate the square-free factorization (Algorithm~8.1) by hand of
$$ a(x) = 1 - 2x + 2 x^3 - x^4 \quad \in {\rm Z}[x].$$
That is, do arithmetic in Z.

\bigskip
\problem Calculate the square-free factorization (Algorithm~8.3) by hand of
$$ b(x) = x^{10} + x^6 + x^5 + x^3 + x^2 + 1 \quad \in {\rm Z}_2[x].$$
That is, do arithmetic in Z$_2$.

\bigskip
\problem (Optional) Factor the polynomials $a(x)$ and $b(x)$ in Z$_2[x]$ using Berlekamp's
algorithm (by hand).  Note that Berlekamp's algorithm works only for 
square-free polynomials.  So we need to do a square-free factorization
first.
$$\eqalign{ a(x) & =  x^{15} - 1, \cr
b(x) & =  x^{10} + x^6 + x^5 + x^3 + x^2 + 1. \cr
}$$

\bye
